\section{Geo-Neus}
Các phương pháp học bề mặt ngầm neural thông qua volume rendering, điển hình là NeuS \cite{Wang2021NeuS} và VolSDF \cite{Yariv2021VolSDF}, đã cho thấy tiềm năng lớn trong việc tái tạo bề mặt 3D từ ảnh đa góc nhìn \cite{Fu2022GeoNeus_source_1}. Tuy nhiên, một thách thức cốt lõi vẫn tồn tại: các phương pháp này chủ yếu dựa vào việc tối ưu hàm mất mát màu sắc (rendering loss) và thiếu các ràng buộc hình học đa góc nhìn (multi-view geometry constraints) một cách tường minh \cite{Fu2022GeoNeus_source_2, Fu2022GeoNeus_source_15, Fu2022GeoNeus_source_33}. Điều này dẫn đến sự không nhất quán về mặt hình học trong bề mặt tái tạo, đặc biệt là ở các vùng ít texture hoặc cấu trúc phức tạp \cite{Fu2022GeoNeus_source_2, Fu2022GeoNeus_source_16}.

Fu và cộng sự \cite{Fu2022GeoNeus} chỉ ra rằng có một "khoảng cách" hay "sai lệch" (gap/bias) cố hữu giữa việc tối ưu tích phân màu sắc trong volume rendering và việc mô hình hóa bề mặt chính xác tại tập mức 0 (zero-level set) của hàm khoảng cách có dấu (SDF) \cite{Fu2022GeoNeus_source_4, Fu2022GeoNeus_source_18, Fu2022GeoNeus_source_19, Fu2022GeoNeus_source_68, Fu2022GeoNeus_source_81, Fu2022GeoNeus_source_87}. Tích phân volume rendering tối ưu trên nhiều điểm dọc tia nhìn, trong khi bề mặt thực sự chỉ được xác định tại một điểm duy nhất (SDF=0). Sự sai lệch này khiến việc chỉ dựa vào mất mát màu sắc để giám sát mạng SDF trở nên không chính xác và khó khăn trong việc tối ưu hình học thực sự \cite{Fu2022GeoNeus_source_24, Fu2022GeoNeus_source_88}.

Để giải quyết vấn đề này, Geo-Neus \cite{Fu2022GeoNeus} được đề xuất với mục tiêu học bề mặt ngầm neural nhất quán về mặt hình học bằng cách đưa vào các ràng buộc hình học đa góc nhìn một cách tường minh.

\subsection{Phương pháp và Mô hình Geo-Neus}

Geo-Neus xây dựng dựa trên nền tảng của NeuS, sử dụng mạng MLP để biểu diễn SDF ($F_\Theta$) và trường bức xạ ($G_\Phi$), cùng với kỹ thuật volume rendering không lệch \cite{Fu2022GeoNeus_source_66, Fu2022GeoNeus_source_78}. Điểm cốt lõi của Geo-Neus là bổ sung hai loại ràng buộc hình học tường minh trực tiếp lên mạng SDF, thay vì chỉ dựa vào giám sát gián tiếp qua mất mát màu sắc \cite{Fu2022GeoNeus_source_3, Fu2022GeoNeus_source_29, Fu2022GeoNeus_source_37, Fu2022GeoNeus_source_90, Fu2022GeoNeus_source_91} (Xem Hình 2 trong \cite{Fu2022GeoNeus}).

\begin{enumerate}
    \item \textbf{Ràng buộc SDF từ Điểm Thưa (Sparse 3D Points - $\mathcal{L}_{SDF}$):}
        \begin{itemize}
            \item \textbf{Nguồn dữ liệu:} Tận dụng các điểm 3D thưa $P$ được tạo ra như một sản phẩm phụ "miễn phí" từ quá trình Structure from Motion (SfM) \cite{Schonberger16SFMRevisited, Snavely2006PhotoTourism} khi ước lượng tham số máy ảnh \cite{Fu2022GeoNeus_source_30, Fu2022GeoNeus_source_94, Fu2022GeoNeus_source_95, Fu2022GeoNeus_source_96}.
            \item \textbf{Giả định:} Các điểm SfM này được coi là nằm trên bề mặt thực của vật thể, do đó giá trị SDF tại các điểm này lý tưởng là bằng 0 ($sdf(p_i) \approx 0$) \cite{Fu2022GeoNeus_source_97, Fu2022GeoNeus_source_98}.
            \item \textbf{Xử lý che khuất:} Với mỗi ảnh đầu vào $I_i$, chỉ sử dụng các điểm SfM $P_i$ nhìn thấy được từ góc nhìn đó để giám sát (dựa trên liên kết với các điểm đặc trưng 2D $X_i$) \cite{Fu2022GeoNeus_source_100, Fu2022GeoNeus_source_101, Fu2022GeoNeus_source_103}.
            \item \textbf{Hàm mất mát:} $\mathcal{L}_{SDF} = \frac{1}{N_i} \sum_{p \in P_i} \| s\hat{d}f(p) \|_1$, với $s\hat{d}f(p)$ là giá trị SDF dự đoán bởi mạng $F_\Theta$. Mất mát này trực tiếp ép buộc mạng SDF dự đoán giá trị 0 tại các điểm đã biết trên bề mặt, đặc biệt hiệu quả ở những vùng có texture phong phú nơi SfM hoạt động tốt \cite{Fu2022GeoNeus_source_104, Fu2022GeoNeus_source_105, Fu2022GeoNeus_source_169}.
        \end{itemize}

    \item \textbf{Ràng buộc Nhất quán Quang trắc Đa góc nhìn (Multi-view Photometric Consistency - $\mathcal{L}_{photo}$):}
        \begin{itemize}
            \item \textbf{Mục tiêu:} Đảm bảo tính nhất quán hình học của bề mặt được biểu diễn bởi SDF=0 trên các góc nhìn khác nhau, đặc biệt hữu ích cho các vùng bề mặt lớn, trơn, ít texture nơi điểm SfM thưa thớt \cite{Fu2022GeoNeus_source_30, Fu2022GeoNeus_source_107, Fu2022GeoNeus_source_108, Fu2022GeoNeus_source_169}.
            \item \textbf{Xác định Điểm Bề mặt trên Tia nhìn:} Thay vì dùng tích phân độ sâu như trong volume rendering (vốn bị bias \cite{Fu2022GeoNeus_source_174, Fu2022GeoNeus_source_178}), Geo-Neus trực tiếp xác định vị trí $t^*$ trên tia nhìn $p(t) = o+tv$ nơi hàm SDF dự đoán $s\hat{d}f(t)$ đổi dấu và bằng 0 (sử dụng nội suy tuyến tính giữa các điểm lấy mẫu rời rạc $t_i, t_{i+1}$ có $s\hat{d}f(t_i) \cdot s\hat{d}f(t_{i+1}) < 0$) \cite{Fu2022GeoNeus_source_31, Fu2022GeoNeus_source_109, Fu2022GeoNeus_source_110, Fu2022GeoNeus_source_114, Fu2022GeoNeus_source_117, Fu2022GeoNeus_source_118}. Điểm giao đầu tiên $t^*$ được chọn để đảm bảo tính nhìn thấy \cite{Fu2022GeoNeus_source_119, Fu2022GeoNeus_source_120}.
            \item \textbf{Áp dụng Nguyên lý MVS:} Với điểm bề mặt $p = o+t^*v$ và pháp tuyến $n = \nabla s\hat{d}f(p)$ tại đó, sử dụng phép biến đổi homography phẳng (plane-induced homography) $H$ để ánh xạ (warp) một vùng ảnh nhỏ (patch) xung quanh điểm chiếu của $p$ trên ảnh tham chiếu $I_r$ sang các ảnh nguồn $I_s$ \cite{Fu2022GeoNeus_source_122, Fu2022GeoNeus_source_125, Fu2022GeoNeus_source_126, Fu2022GeoNeus_source_127, Hartley2004MVG}.
            \item \textbf{Đo lường sự nhất quán:} Tính toán sự tương đồng quang trắc giữa patch trên ảnh tham chiếu và các patch được ánh xạ trên ảnh nguồn bằng cách sử dụng Normalized Cross-Correlation (NCC) trên ảnh xám \cite{Fu2022GeoNeus_source_128, Fu2022GeoNeus_source_130, Furukawa2010MVS, Galliani2015Gipuma}.
            \item \textbf{Hàm mất mát:} $\mathcal{L}_{photo}$ được định nghĩa dựa trên $1 - NCC$ của các patch tương ứng, thường lấy trung bình của một số (ví dụ: 4) giá trị NCC tốt nhất để xử lý che khuất \cite{Fu2022GeoNeus_source_132}. Mất mát này khuyến khích hình học (vị trí $p$ và pháp tuyến $n$) phải nhất quán giữa các góc nhìn khác nhau \cite{Fu2022GeoNeus_source_133}.
        \end{itemize}

    \item \textbf{Hàm Mất mát Tổng thể:} Kết hợp các thành phần:
        $$\mathcal{L} = \mathcal{L}_{color} + \alpha\mathcal{L}_{reg} + \beta\mathcal{L}_{SDF} + \gamma\mathcal{L}_{photo}$$
        trong đó $\mathcal{L}_{color}$ là mất mát màu sắc (L1 loss), $\mathcal{L}_{reg}$ là ràng buộc Eikonal ($\|\nabla s\hat{d}f\| \approx 1$) \cite{Gropp2020IGR}, $\mathcal{L}_{SDF}$ và $\mathcal{L}_{photo}$ là các mất mát hình học mới. Các hệ số $\alpha, \beta, \gamma$ cân bằng các thành phần \cite{Fu2022GeoNeus_source_134}.
\end{enumerate}

Kiến trúc mạng MLP cho SDF và màu sắc tương tự như trong NeuS \cite{Wang2021NeuS} và IDR \cite{Yariv2020IDR}, sử dụng positional encoding \cite{Mildenhall2020NeRF} và khởi tạo hình học \cite{Atzmon2020SAL} \cite{Fu2022GeoNeus_source_145, Fu2022GeoNeus_source_146, Fu2022GeoNeus_source_147}.

\subsection{Phân tích Ưu nhược điểm}

\subsubsection{Ưu điểm}
\begin{itemize}
    \item \textbf{Chất lượng Hình học Vượt trội:} Nhờ các ràng buộc hình học đa góc nhìn tường minh ($\mathcal{L}_{SDF}$, $\mathcal{L}_{photo}$), Geo-Neus tạo ra các bề mặt có độ chính xác và nhất quán hình học cao hơn đáng kể so với các phương pháp chỉ dựa vào giám sát màu sắc như NeuS, VolSDF \cite{Fu2022GeoNeus_source_2, Fu2022GeoNeus_source_3, Fu2022GeoNeus_source_7, Fu2022GeoNeus_source_32, Fu2022GeoNeus_source_34, Fu2022GeoNeus_source_133, Fu2022GeoNeus_source_154, Fu2022GeoNeus_source_155, Fu2022GeoNeus_source_187}.
    \item \textbf{Giảm Sai lệch (Bias) của Volume Rendering:} Việc tối ưu trực tiếp trên bề mặt SDF=0 thông qua các ràng buộc hình học giúp giảm thiểu tác động của sai lệch vốn có trong phương pháp tích phân màu của volume rendering \cite{Fu2022GeoNeus_source_4, Fu2022GeoNeus_source_6, Fu2022GeoNeus_source_31, Fu2022GeoNeus_source_38, Fu2022GeoNeus_source_87, Fu2022GeoNeus_source_178}.
    \item \textbf{Tái tạo Tốt Cấu trúc Phức tạp và Vùng Trơn:} Sự kết hợp của $\mathcal{L}_{SDF}$ (tập trung vào vùng có texture, cấu trúc mỏng) và $\mathcal{L}_{photo}$ (tập trung vào vùng trơn, ít texture) giúp Geo-Neus tái tạo tốt cả hai loại bề mặt này \cite{Fu2022GeoNeus_source_7, Fu2022GeoNeus_source_34, Fu2022GeoNeus_source_40, Fu2022GeoNeus_source_156, Fu2022GeoNeus_source_169, Fu2022GeoNeus_source_170, Fu2022GeoNeus_source_186}.
    \item \textbf{Hội tụ Nhanh hơn:} Giám sát hình học tường minh giúp mạng hội tụ nhanh hơn và ổn định hơn so với NeuS \cite{Fu2022GeoNeus_source_105, Fu2022GeoNeus_source_180, Fu2022GeoNeus_source_182}.
\end{itemize}

\subsubsection{Nhược điểm và Hạn chế}
\begin{itemize}
    \item \textbf{Chi phí Tính toán và Thời gian Huấn luyện:} Mặc dù hội tụ nhanh hơn về số vòng lặp, việc tính toán thêm các ràng buộc hình học (đặc biệt là $\mathcal{L}_{photo}$ với warping và NCC) làm tăng chi phí tính toán trên mỗi vòng lặp. Thời gian huấn luyện tổng thể vẫn còn đáng kể (đề cập khoảng 16 giờ), tương đương NeuS và chậm hơn nhiều so với các phương pháp tăng tốc như Instant-NGP hay NeuS2 \cite{Fu2022GeoNeus_source_149, Fu2022GeoNeus_source_188}.
    \item \textbf{Phụ thuộc vào SfM:} Chất lượng và độ phủ của các điểm 3D thưa từ SfM có thể ảnh hưởng đến hiệu quả của $\mathcal{L}_{SDF}$. Các vùng hoàn toàn không có texture sẽ không nhận được giám sát từ loại ràng buộc này.
    \item \textbf{Độ phức tạp:} Mô hình và quy trình huấn luyện phức tạp hơn do có thêm các thành phần xử lý hình học (SfM, MVS).
    \item \textbf{Không đề cập đến Cảnh động:} Bài báo chỉ tập trung vào tái tạo cảnh tĩnh.
\end{itemize}

Tóm lại, Geo-Neus là một cải tiến quan trọng cho các phương pháp học bề mặt ngầm neural bằng cách tích hợp các ràng buộc hình học đa góc nhìn tường minh, giúp khắc phục sai lệch của volume rendering và cải thiện đáng kể chất lượng, độ nhất quán của bề mặt tái tạo, dù vẫn còn hạn chế về tốc độ tính toán.
